\scanpage{46}
\mspage{5}
% L1–6
une base s'ils sont générateurs. Inversement, si ils font une base, $\exists!$ forme bil. sym. satisfaisant les conditions équivalentes (i), (i'). On aura alors les formules (12) à (16) de (I). On suppose dorénavant qu'on est dans cette situation.

% L7–13
\unsure{§} Considérons les \unsure{27|17} éléments $\unsure{\zeta|\xi}_i,\zeta'_i$ \inserted{$(1\le i\le6)$}, $\zeta_{ij}$ $(1\le i\ne j\le6)$, qui dans la base $(\eta,(\zeta_i))$ s'écrivent
\[
\begin{aligned}
\zeta_i&=\zeta_i\\
\zeta'_i&=2\eta-\sum_{\alpha\ne j}\zeta_i\\
\zeta_{ij}&=\eta-\zeta_i-\zeta_j\quad\inserted{(i\ne j)}
\end{aligned}
\]
\margin{\uline{NB} \deleted{les} $\delta$ engendrent $\widehat E$ (il suffit $=$ des $\zeta_i$ et d'\uline{un} $\zeta_{ij}$)}
donc sont tous distincts. \inserted{Soit \deleted{\unsure{$\Delta$ leur}} $\Delta_{\unsure{0|X}}$ leur ens.} On met dessus une structure de graphe en notant \deleted{\unsure{convenant}} que si $\delta,\delta'\in\Delta_{\unsure{x|k|X}}$, \inserted{$\delta\ne\delta'$}, on a $\delta\cdot\delta'=0$ ou 1, \inserted{(NB $\delta^2=\delta'{}^2=-1$)} et \unsure{en considérant|on considère} \deleted{$\delta,\delta'$ liés sss \unsure{à|$\delta$}}
\[
\delta,\delta'\ \inserted{\in\Delta_x,\ \delta\ne\delta'}\ \text{liés ssi}\qquad\delta\delta'=1\qquad\text{i.e. }\delta\delta'\ne0.\tag*{(j)}
\]
% L14–21
Donc $\zeta_i$ est lié aux él. $\zeta'_j$ $(j\ne i)$ et $\zeta_{ij}$ $(j\ne i)$ \unsure{exclt|exactement} (soit 10 él.), $\zeta'_i$ est lié aux él. $\zeta_j$ $(j\ne i)$ et $\zeta_{ij}$ $(j\ne i)$ \unsure{exclt|exactement} (soit encore 10 éléments) et $\zeta_{ij}$ \inserted{$(i\ne j)$} est lié à $\zeta_i,\zeta_j,$ \inserted{$\zeta'_i,\zeta'_j,$} et à $\zeta_{k\ell}$ $(k\ne i,j,\ell\ne i,j,k\ne\ell)$ en nombre $\binom42=6$, donc en tout encore 10 éléments. \deleted{Enfin,} Les \deleted{couples} \inserted{paires} de sommets liés sont donc
% archive sheet number: 45
